\section{Local bounds, planar filling, and positive temporal deficit}
\label{sec:local}

We work at the fixed critical parameter point from
Lemma~\ref{lem:critical-point}.  The arguments in this section first
apply to a hybrid equality from Lemma~\ref{lem:hybrid-equality}, with
\[
 0<h_0<H\leq h,\qquad
 m\eqexp\Delta_{h_0}N^{H+\Pi_F(1)}.
\]
They also apply to an exact equality, when $H=h_0=h>0$.
Our goal is to obtain an exact equality with positive total deficit
and an initial interval of zero deficit.  Zero deficit means full
temporal branching in exponent; the actual filling of time blocks
will be a separate operation.

The argument has two geometric stages.  Multilinear Kakeya and the
planar Furstenberg estimate control the support cost of filling short
time blocks, forcing an initial flat interval in every equality
profile.  The weighted set estimate then captures plate packets that attain
the local mass bound on those flat scales.  Anisotropic rescaling of aligned plates rules
out arbitrarily small total deficit for hybrid equalities near $h$.
Its strict gain comes from $H-h_0>0$; when $H=h_0$, the same
rescaling will produce new exact equalities on shorter time intervals
in Section~\ref{sec:stationary}.

We abbreviate $\Pi_F$ to $\Pi$ and set
\[
 p_*=p-\tfrac12>\tfrac32,\qquad
 x_H=d-1-H,\quad y_H=1+H.
\]
In particular $x_H+y_H=d$, $x_H\geq0$, and
\begin{equation}
 \Pi(u)\geq p_*F(u),\qquad
 \Pi(1)\geq p\alpha-q.
 \label{loc:cost-lower}
\end{equation}
The first inequality follows from
$\mathcal D(q,e)\leq e/2$, and the second from
$\mathcal D(q,e)\leq q$.

All clustering and multiplicity estimates use the inherited weights
$\omega_i$.  The probability measures in the planar Frostman argument
will have their normalizing denominators displayed.  An extensive
regular refinement preserves $F$, $m$, and $\Delta_{h_0}$ in exponent,
by Lemma~\ref{lem:hybrid-equality}; its sharp-bound comparison also
applies when $H=h_0=h$.  We use this stability after each of the
finitely many restrictions at a fixed stage of the slow diagonal.

In local geometric arguments, lengths such as $a,b,K$ are in original
fine-cell units.  Thus a temporal block of $K$ bins has actual duration
$K/N$, and a transverse width $a$ means actual width $a/N$.
Containment of an index in a plank over a block always means containment
of its entire affine segment over that block.

\subsection{The local mass bound and the isotropic scale cut}

An index is active on a temporal block if that block contains one
of its retained marked bins.  A packet on the block is a set of
active indices assigned to a plank containing their entire affine
segments over that block.  Its mass is the sum of their inherited
weights.  Disjoint packet assignments mean disjoint assigned index
sets on each block; their containing planks may overlap.

\begin{proposition}[Local bound]
\label{prop:local-bound}
After an extensive regular restriction, the mass of active original
indices contained in any plank of widths $1\leq a\leq b\leq K$ over
a dyadic temporal block of $K=N^\kappa$ bins satisfies
\begin{equation}
 n_{I,W}\leexp T(\kappa,a,b),\qquad
 T(\kappa,a,b)=
 mN^{-\Pi(\kappa)}K^{-H}a^{d-1-H}b^{1+H}.
 \label{loc:mass-bound}
\end{equation}
The bounds hold simultaneously at every fixed limiting scale and
shape, under the finite-test and slow-diagonal convention of
Section~\ref{sec:critical}, and are inherited by further restrictions.

Suppose conversely that, on an extensive graph, disjoint indexed-block
assignments to planks of common widths $a,b$ have masses at least
$T(\kappa,a,b)$ in exponent.  Set
\begin{equation}
 \mathfrak a=\log_N a,
 \qquad \mathfrak u=\log_N(K/b),
 \qquad \chi=\mathfrak a+\mathfrak u.
 \label{loc:cut-loss}
\end{equation}
If $\chi<1$, the isotropic scale cut described below gives an ensemble
at resolution $N'=Nb/(Ka)=N^{1-\chi}$ with matching sharp upper and
lower bounds.  At $H=h_0=h$ this is an ensemble of exact critical
equalities, from which a realizing representative sequence can be
selected.  Its temporal profile is
\begin{equation}
 G(v)=
 \begin{cases}
  0,&0\leq v\leq(\kappa-\chi)/(1-\chi),\\[2pt]
  \displaystyle
  \frac{F(\chi+(1-\chi)v)-F(\kappa)}{1-\chi},
     &(\kappa-\chi)/(1-\chi)\leq v\leq1.
 \end{cases}
 \label{loc:cut-profile}
\end{equation}
\end{proposition}

\begin{proof}
\emph{The rescaling and its cell count.}
We first describe the operation that determines the threshold $T$.
Suppose active indices on selected $K$-blocks have been assigned
to packets of widths $a,b$.  Suppose also that the indices of each packet lie in one cell
of an isotropic grid in matrix space, of side $b/K$; call this cell
its full-time parent.  Subtract the parent center matrix and divide
space by $b/K$, leaving time unchanged.  Use resolution
\begin{equation}
 N'=\frac{Nb}{Ka}.
 \label{loc:isotropic-resolution}
\end{equation}
The new matrices lie in a bounded patch.  A new spatial fine cell
corresponds to $a$ old spatial fine cells, and an original $K$-block
contains $b/a$ new time bins.

Fill each selected block with those new bins, keeping all selected
blocks of one original index together in its full-time parent.
A packet now has transverse widths $1,b/a$.  It occupies at most
$C(b/a)^2$ space--time cells: at most $Cb/a$ spatial cells per bin,
over at most $Cb/a$ bins.  If its inherited index mass is $M$, its
incidence mass is at least $cMb/a$.  Thus a collection of disjoint
indexed-block assignments, each of mass at least $M$, has aggregate
raw multiplicity at least a constant times $Ma/b$.  We count the
supports in different parents separately and bound each union by the
sum of its packet supports.  Overlaps within a parent only improve
this lower bound.

This calculation explains the factor $a/b$ in the comparison below.
To use the sharp upper bound we must also preserve the original time
tree above scale $K$.  We now construct the extensive parent ensemble
for which that profile comparison is valid.

\paragraph{Capturing and routing an excessive family.}
Fix a proposed excess $N^\varepsilon$ and a finite scale and shape
test.  On each temporal block, greedily assign all remaining active
indices in a plank with mass greater than
$N^\varepsilon T(\kappa,a,b)$, and remove their events on that
block.  Continue until no such plank remains.  The plank tests may
be meshed and thickened by a fixed factor, so the procedure is finite.
Its assignments are disjoint in original indices on each block.
Suppose the removed events form an extensive graph.  Pigeonholing
dyadic widths costs only a subpower, so retain a common pair $a,b$.

Two matrices whose segments stay in the same local plank differ in
spatial slope by $O(b/K)$: evaluate the difference at the endpoints
of the block and divide its $O(b/N)$ transverse diameter by the
duration $K/N$.  Their intercepts likewise differ by $O(b/K)$,
since all original times belong to a bounded interval.  Each packet
therefore meets only boundedly many full-time isotropic matrix-grid
parents of side $b/K$.  Partition its indices among those parents.
For a fixed width type each original matrix has one parent, so an
index is not duplicated once for each selected block.

Before this subdivision, every assigned index-block contributes
$KN^{-F(\kappa)+o(1)}$ events.  Include the assignment cells and
their parent subdivisions among the regularization tests.  The
inherited-weight cleaning of Lemma~\ref{lem:regularization} and
Corollary~\ref{cor:active-weight-cleaning} preserves the threshold $N^{\varepsilon-o(1)}T$ on surviving
packets: delete active packets retaining less than an inverse
subpower fraction of their original incidence, and retain
nondeficient parent subdivisions.  Upper branching is inherited
on each original index and starting block.  Total incidence
retention and bounded routing then preserve the whole profile
$F$ on the parent ensemble.

Apply the rescaling and filling just described.  Its cell count
gives aggregate raw multiplicity at least
\begin{equation}
 N^{\varepsilon-o(1)}T(\kappa,a,b)\frac{a}{b}.
 \label{loc:isotropic-lower}
\end{equation}

\paragraph{The sharp upper bound.}
A new full-time plank of widths $A,B$ pulls back to one of widths
$aA,aB$ in original fine units.  These are legitimate original tests:
$1\leq aA\leq aB\leq aN'=Nb/K\leq N$.  Since the shape exponents
for $\Delta_{h_0}$ sum to $d$, each parent's raw parameter obeys
\begin{equation}
 \Delta_{h_0,\mathrm{new}}\leexp a^d\Delta_{h_0}.
 \label{loc:isotropic-parameter}
\end{equation}
The old occupied time tree above scale $K$ is unchanged; below it we
have grafted full time intervals of $b/a$ bins.  The profile is exactly
\eqref{loc:cut-profile} in normalized logarithmic coordinates.
In particular its temporal cost, expressed in original $\log N$
units, is $\Pi(1)-\Pi(\kappa)$.

The sharp hybrid upper bound in each parent is consequently at most
\begin{align}
 a^d\Delta_{h_0}(N')^H
       N^{\Pi(1)-\Pi(\kappa)}
 &\eqexp mN^{-\Pi(\kappa)}K^{-H}a^{d-H}b^H \notag\\
 &=T(\kappa,a,b)\frac{a}{b}.
 \label{loc:isotropic-upper}
\end{align}
The same upper bound holds uniformly for the ensemble.  Otherwise a
sequence of failing parent representatives would contradict the
definition of the sharp exponent after regularizing to their common
limiting profile.  This contradicts
\eqref{loc:isotropic-lower}.

If $N'=N^{o(1)}$, then $\chi\to1$.  Because
$\chi\leq\kappa\leq1$, one has $\kappa\to1$ and
$\Pi(1)-\Pi(\kappa)=o(1)$.  The full-box mass bound in a parent,
at most $a^d\Delta_{h_0}(N')^d$, supplies the same exponent upper
bound as \eqref{loc:isotropic-upper}.  Thus this endpoint also cannot
carry the proposed fixed excess.

It follows that removing excessive assignments costs no extensive
piece.  Perform these removals successively for a finite mesh of
scales and shapes, regularizing after each removal.  The observation
at the start of the section preserves equality for the next test.
Nearby exponents are controlled by enlarging widths and extending to
nearby nested temporal blocks; endpoint extrapolation costs only the
prescribed mesh error.  Let the mesh and error tolerances tend to
zero in the stated order.  This proves \eqref{loc:mass-bound}.

For the converse assertion, start with the extensive saturated
assignments in the statement and perform the same routing, rescaling,
and filling, with packet mass at least $T$ in exponent and no fixed
excess.  The lower bound is now $Ta/b$ and the upper bound is exactly
the same expression.  When $\chi<1$, the derived resolution grows with a
positive exponent and the profile in \eqref{loc:cut-profile} is
well-defined.  Equality of the ensemble bounds provides either the
ensemble formulation or realizing representatives.  At $H=h_0=h$
these are equalities for the original critical problem.
\end{proof}

\subsection{Multilinear slack and a common plane on each short block}

We first isolate two geometric consequences of the local bound.  They
will be used both in the flat-delay argument and in the later alignment
argument.

\begin{lemma}[Slack and cap mass]
\label{loc:slack-caps}
Every equality under consideration satisfies
\begin{equation}
 N^{H+(p-2)\alpha}\leexp\frac{k^2m}{n}.
 \label{loc:multilinear-slack}
\end{equation}
At a fine incidence, the raw mass of active lines whose directions
lie in any cap of radius $C/L$, $1\leq L\leq N$, is at most
\begin{equation}
 N^{o(1)}mL^{-H}N^{-\Pi(\log_N L)}.
 \label{loc:cap-bound}
\end{equation}
The constant $C$ is fixed; the exponent error is uniform under the
finite-test convention.
\end{lemma}

\begin{proof}
The full bounded patch is covered by boundedly many planks with both
widths comparable to $N$, so $n\leexp\Delta_{h_0}N^d$.  Since
$k\eqexp N^{1-\alpha}$ and
$m\eqexp\Delta_{h_0}N^{H+\Pi(1)}$, the ratio in
\eqref{loc:multilinear-slack} is at least
$N^{q+H+\Pi(1)-2\alpha}$ in exponent.  Inequality
\eqref{loc:cost-lower} gives the claim.

Lines meeting a fixed fine cell with slope diameter $O(1/L)$ stay
within a bounded-width plank about one of them throughout the
dyadic $L$-block containing that cell.  They are active on that block.
Apply \eqref{loc:mass-bound} with $a,b=O(1)$, using bounded coverings
to handle fixed enlargements.  Bounded slope charts identify spatial
slope diameter and spherical direction diameter up to constants.
This proves \eqref{loc:cap-bound}.
\end{proof}

The concentration of directions near a plane is the planiness
principle familiar from \citet{KLT2000}; see also the related
multilinear consequence in \citet[Corollary~1.4]{Guth2010}.
For filling time blocks, we need the direction plane to stay fixed
along each reference index throughout the block.  We obtain this control
by applying the multilinear estimate on balls of the block scale;
the positive slack makes the incidence in balls with three rich
transverse directions negligible.

\begin{lemma}[Common short-block plane]
\label{loc:short-block-plane}
There is $\kappa_{\mathrm{ml}}>0$, depending only on the positive
slack in \eqref{loc:multilinear-slack} and the multilinear input,
with the following property.  For each fixed
$0<\kappa<\kappa_{\mathrm{ml}}$, put $K=N^\kappa$.  On an extensive
regular central incidence graph, each active reference index-block
can be assigned one direction plane such that:
\begin{enumerate}
 \item the same plane is used at all of its retained times;
 \item at each such time there is original neighboring incidence mass
 at least $mN^{-o(1)}$, in directions within $O(K^{-1})$ of that plane;
 \item after discarding a negligible part of this neighboring mass,
 the neighbor direction is separated from the reference direction
 by at least $N^{-o(1)}$.
\end{enumerate}
All reference and neighbor segments over the block lie in a common
plate of dimensions $C\times CK\times CK$ in fine units.  The
neighboring indices retain their original marked-time shadings and
their inherited raw cap bounds.
\end{lemma}

\begin{proof}
Choose a lattice cover by balls of radius a sufficiently large fixed
multiple of $K$ in fine space-time units.  Its overlap is bounded.
The radius is large enough that one of the balls assigned to a
reference index-block contains its whole segment and the whole
$K$-block segment of every bounded-slope line meeting it at a fine
cell.  Only boundedly many choices are needed for any index-block.

By \eqref{loc:mass-bound}, the active mass near such a ball is at
most $mK^{C_0}$, for an absolute constant $C_0$: cover its time
extent by boundedly many $K$-blocks and its transverse extent by
boundedly many planks with $a,b\asymp K$.  Multiplying by the at
most $CK$ events of each active line gives an incidence upper bound
$mK^{C_0+1}$ for the ball.

Divide directions into caps of radius $K^{-20}$ and call a cap rich
at the ball if its active mass there is at least $mK^{-100}$.
There are $O(K^{40})$ caps.  A ball is bad if three rich caps have
direction determinant at least $K^{-10}$.  Apply
Theorem~\ref{thm:multilinear-input} to the unshaded tubes at spatial
resolution $K/N$, summing over the possible cap triples.  The
transversality loss, the cap count, and the thresholds are fixed
powers of $K$.  Therefore, for a fixed constant $C_1$ and any fixed
$\varepsilon>0$,
\begin{equation}
 \#\{\text{bad balls}\}
 \lesssim_\varepsilon N^\varepsilon K^{C_1}(n/m)^{3/2}.
 \label{loc:bad-balls}
\end{equation}
Indeed, the multilinear summand at a tested ball is at least
$(mK^{-100})^{3/2}$, and the product of the three full tube masses
is at most $n^{3/2}$.  Bounded ball-to-cube comparisons do not change
the estimate.

The fraction of all weighted incidence in bad balls is consequently
at most
\[
 N^\varepsilon K^{C_2}
       \left(\frac{n}{k^2m}\right)^{1/2}.
\]
Choose first $\kappa_{\mathrm{ml}}$ small enough that the $K^{C_2}$
factor uses less than one quarter of the positive slack exponent,
and then choose $\varepsilon$ smaller still.  By
\eqref{loc:multilinear-slack}, bad balls carry a power-small
incidence fraction.  The same estimate applies to the bounded
enlargements and choices of balls just used.

At a remaining ball, any three rich directions have determinant
$O(K^{-10})$.  Choose a maximally separated pair.  If its separation
is at most $K^{-1}$, all rich directions are within $O(K^{-1})$ of
any plane containing one of them.  Otherwise the norm of the pair's
cross product is bounded below by $c/K$, and dividing the determinant
bound by this norm places every rich direction within $O(K^{-9})$
of their spanning plane.  In either case there is a single plane
with the claimed $O(K^{-1})$ tolerance.  Its normal has spatial
component bounded below: it is approximately perpendicular to a
direction of the form $(1,u)$ with bounded $u$.

The total mass in light caps at the ball is at most $CmK^{-60}$.
Cells whose original multiplicity is less than $mN^{-\varepsilon_1}$
carry at most $N^{-\varepsilon_1}$ of all incidence, because their
number is at most $|E|$.  Let $\varepsilon_1$ decrease slowly to zero.
Remove central incidences requiring bad balls and central incidences
in the light caps of their assigned ball.  At each relevant fine
cell the light-cap bound is negligible compared with
$mN^{-o(1)}$; the bounded number of possible adjacent balls does not
alter that conclusion.  Choose the ball once for each retained
reference index-block, and keep this choice fixed at all its times.
Regularize this central graph, preserving $F$.

For neighbors at a retained time use the original rich incidences at
that fine cell.  Their total mass is at least $mN^{-o(1)}$.
The cap bound \eqref{loc:cap-bound} with a radius $N^{-\gamma}$
shows that neighbors within that angle of the reference have mass
at most $mN^{-H\gamma+o(1)}$.  First take any fixed small
$\gamma>0$, then diagonalize $\gamma\downarrow0$ after the previous
errors are smaller.  This leaves separation at least $N^{-o(1)}$
and does not use an estimate at an uncontrolled shrinking scale.

Finally, all these directions are $O(K^{-1})$ from the one chosen
plane, and each neighbor meets the reference within bounded fine
distance.  Over a block of length $K$, their displacement normal
to the affine plane through the reference remains bounded.  Their
tangential and temporal displacements are $O(K)$.  A plate of the
stated dimensions contains them all, including the filled reference
segment.  Neighbor shadings have not been restricted in this step,
so their original profile and cap upper bounds remain available.
\end{proof}

\subsection{Planar filling and the flat delay}

The common short-block plane supplies a plate containing each full
reference segment.  We will use its neighboring shadings to show
that the original union has substantial density in this plate, and
then fill the reference segment.  Two estimates control the support
added by this operation.  The planar Furstenberg theorem supplies
the plate density; the following maximal estimate converts that
density into a bound for the union of all filled segments.

\begin{lemma}[Plate Nikodym bound]
\label{lem:plate-nikodym}
For $K\geq2$, let $\mathcal N_K$ be the maximal averaging operator
over rectangular plates of dimensions $1\times K\times K$, with
arbitrary orientations and translations, whose averaging plate
contains the query point.  Fixed changes in these dimensions are
allowed.  Then
\begin{equation}
 \|\mathcal N_K f\|_{L^2(\R^3)}^2
       \lesssim(1+\log K)\|f\|_{L^2(\R^3)}^2.
 \label{loc:nikodym-l2}
\end{equation}
Consequently, if every point of a measurable set $E^+$ belongs to
one such plate in which $E$ has density at least $\rho>0$, then
\begin{equation}
 |E^+|\lesssim \rho^{-2}(1+\log K)|E|.
 \label{loc:nikodym-density}
\end{equation}
The same statement holds for fine-cell counts after bounded
thickening.
\end{lemma}

\begin{proof}
Take any measurable linearized choice $P_x$ of an averaging plate
containing $x$, and put
\[
 Af(x)=\frac1{|P_x|}\int_{P_x}f(z)\,dz.
\]
The symmetric kernel of $AA^*$ is
\[
 \mathcal K(x,y)=\frac{|P_x\cap P_y|}{|P_x||P_y|}.
\]
Let $\vartheta\in[0,\pi/2]$ be the angle between the unoriented
central-plane normals of the two plates.  Intersecting two unit
slabs gives transverse cross-sectional area $O(1/\sin\vartheta)$,
while the available extent along their intersection line is $O(K)$.
The trivial overlap bound is $O(K^2)$.  Thus
\begin{equation}
 |P_x\cap P_y|\lesssim
       \min(K^2,K/\vartheta)
       \lesssim\frac{K^2}{1+K\vartheta}.
 \label{loc:plate-overlap}
\end{equation}

Fix $P_x$ and let $\mathcal P_x$ be its central plane.  If the plates
intersect, choose a point in their intersection.  Since $y\in P_y$,
the normal displacement from that point to $y$, measured relative
to $\mathcal P_x$, is at most $C(1+K\vartheta)$.  It follows that
\[
 \dist(y,\mathcal P_x)\leq C(1+K\vartheta).
\]
Also $|x-y|\leq CK$, since both plates contain their respective
query points and have a common point.  Because $|P_x|,|P_y|\asymp K^2$,
\eqref{loc:plate-overlap} gives
\begin{equation}
 0\leq\mathcal K(x,y)
 \lesssim
 \frac{\one_{\{|x-y|\leq CK\}}}
      {K^2(1+\dist(y,\mathcal P_x))}.
 \label{loc:schur-kernel}
\end{equation}
In coordinates tangent and normal to $\mathcal P_x$, the tangential
area in this ball is $O(K^2)$.  Integrating the normal coordinate
therefore proves
\[
 \sup_x\int\mathcal K(x,y)\,dy\lesssim1+\log K.
\]
Symmetry gives the same column bound.  Schur's test yields
$\|AA^*\|_{2\to2}\lesssim1+\log K$, proving the claimed squared
norm bound uniformly in the linearization.

For a finite family of plate parameters, choose a maximizing plate
measurably at each point.  The uniform bound passes to the supremum
over that family.  A countable dense family suffices, by continuity
of plate integrals under translation and rotation for locally
integrable functions, and monotone convergence then gives
\eqref{loc:nikodym-l2}.  Apply it to $f=\one_E$ and use
$\mathcal N_K\one_E\geq\rho$ on $E^+$ to obtain
\eqref{loc:nikodym-density}.  Replacing counted cells by their unit
cubes and allowing fixed plate enlargements proves the discrete
version.
\end{proof}

To obtain the plate density, sample a marked reference time and then
a separated neighboring line at that time.  If the time distribution
has Frostman exponent $s$ and the conditional direction distribution
has exponent $a_{\mathrm{ang}}$, their joint line-parameter law has
exponent $t=s+a_{\mathrm{ang}}$.  The planar theorem contributes
$\mu=\min\{1,(s+t)/2,t\}$ to the union size.  The next lemma
compares the resulting support cost of filling with the temporal
cost $\Pi(\kappa)$ that filling removes.

\begin{lemma}[Planar-prefix extension]
\label{lem:prefix-extension}
Let $0<\kappa<\kappa_{\mathrm{ml}}$ be fixed as in
Lemma~\ref{loc:short-block-plane}, and put $K=N^\kappa$.  Suppose
there are constants $s\in(0,1]$ and $\gamma\geq0$ such that
\begin{equation}
 F(\kappa)-F(u)\leq(1-s)(\kappa-u),
 \qquad F(u)\geq\gamma u
 \quad(0\leq u\leq\kappa).
 \label{loc:prefix-assumptions}
\end{equation}
Set
\begin{equation}
 a_{\mathrm{ang}}=H+p_*\gamma,
 \qquad t=s+a_{\mathrm{ang}},
 \qquad \mu=\min\{1,(s+t)/2,t\}.
 \label{loc:prefix-exponents}
\end{equation}
Then $a_{\mathrm{ang}}\leq1$, so $0<t\leq2$.  Every equality
configuration under consideration must satisfy
\begin{equation}
 \Pi(\kappa)\leq F(\kappa)+2\kappa(1-\mu).
 \label{loc:prefix-necessary}
\end{equation}
In particular, a strict reverse inequality yields a contradiction.
There is no requirement that $\gamma$ or $F(\kappa)/\kappa$ have a
positive lower bound.  The small-scale range is uniform over
equalities with $H$ bounded below by a fixed positive constant.
\end{lemma}

\begin{proof}
\emph{The planar neighbor measure.}
Use the central graph and the fixed block planes of
Lemma~\ref{loc:short-block-plane}.  Choose a retained reference
index-block.  Normalize its block length to one and project
orthogonally onto the corresponding affine plane.  The projection
changes its relevant unit directions by an additive $O(K^{-1})$.
The reference--neighbor separation is $N^{-o(1)}$, which is much
larger than $K^{-1}$ for this fixed $\kappa>0$; that separation
is therefore preserved up to fixed factors.  Also, projected
direction diameter $O(r)$ implies original direction diameter
$O(r+K^{-1})=O(r)$ whenever $r\geq K^{-1}$.

We use physical time, rather than distance along the reference,
as the graph coordinate on this plane.  After a translation let
its unit normal be $\nu=(\nu_0,\nu_{\mathrm{sp}})$, where
$|\nu_{\mathrm{sp}}|\geq c>0$.  Put
$n_{\mathrm{sp}}=\nu_{\mathrm{sp}}/|\nu_{\mathrm{sp}}|$ and
choose a unit spatial vector $e_{\mathrm{sp}}$ perpendicular to it.
The coordinates $(\sigma,y)=(t,e_{\mathrm{sp}}\cdot x)$ on the
plane have inverse
\[
 (t,x)=\left(\sigma,
 y e_{\mathrm{sp}}-
 \frac{\nu_0}{|\nu_{\mathrm{sp}}|}\sigma n_{\mathrm{sp}}\right),
\]
so they are uniformly bi-Lipschitz.  For a relevant original
direction $(1,u)$, the projected time component is
$1-\nu_0\nu\cdot(1,u)=1+O(K^{-1})$, while its projected
$y$ component is $e_{\mathrm{sp}}\cdot u$.  All projected lines
are consequently bounded-slope graphs in these coordinates.
This also preserves lengths, time resolution, and shading counts
up to fixed factors.

First, the marked-time probability distribution on the reference
block is $s$-Frostman at resolution $K^{-1}$.  To check this, an
interval of relative length $r\in[K^{-1},1]$ is covered by boundedly
many dyadic blocks of $L\asymp rK$ bins.  Put $u=\log_N L$.
The relative number of reference events it contains is at most
\begin{equation}
 N^{o(1)}\frac{LN^{-F(u)}}{KN^{-F(\kappa)}}
 \leq N^{o(1)}r^s,
 \label{loc:time-frostman}
\end{equation}
by \eqref{loc:prefix-assumptions}.  The same estimate holds for
the original marked times of every neighboring line on this block.
The errors are subpowers of $K$ because $\kappa>0$ is fixed.

Construct a probability measure on planar neighbor lines as follows.
Choose a marked reference time with the preceding distribution,
then choose one of its coplanar separated neighbors with probability
proportional to its original raw weight at that fine cell.  The
normalizing neighboring mass is at least $mN^{-o(1)}$.
For any cap of radius $r\in[K^{-1},1]$, use
\eqref{loc:cap-bound} with $L\asymp1/r$ and
\eqref{loc:cost-lower} to obtain the conditional probability bound
\begin{equation}
 \Prob(\text{neighbor direction in that cap}\mid\text{time})
 \leq N^{o(1)}r^H
       N^{-p_*F(\log_N(1/r))}
 \leq N^{o(1)}r^{a_{\mathrm{ang}}}.
 \label{loc:conditional-cap}
\end{equation}
The coplanar directions fit in $O(K)$ caps of radius $K^{-1}$.
Their normalized mass is one, so summing
\eqref{loc:conditional-cap} at this radius gives
$1\leq N^{o(1)}K^{1-a_{\mathrm{ang}}}$.  Hence
$a_{\mathrm{ang}}\leq1$ in the limit.

We claim that the resulting planar line-parameter probability measure
is $t$-Frostman.  A line-parameter ball of radius $r$ restricts the
projected slope to an interval of size $O(r)$, and therefore the
original three-dimensional direction to a cap of size $O(r)$:
the lost normal component is only $O(K^{-1})\leq O(r)$.  It also
restricts the encounter time with the reference to an interval of
length $N^{o(1)}r$.  For the latter assertion, subtract the
projected reference's affine graph by a bounded shear in the
physical-time coordinates just defined, and write a projected
neighbor as $y=c+v\sigma$.  The neighbor-reference angle
is at least $N^{-o(1)}$, so $|v|\geq N^{-o(1)}$.  The encounter
condition is $|c+v\sigma|\lesssim K^{-1}$.  Varying $(c,v)$ in an
$r$-ball changes $-c/v$ by at most $N^{o(1)}r$, provided the ball
is smaller than a fixed fraction of $|v|$.  For larger balls the
same assertion follows from the trivial unit bound on the time
interval.  The fine encounter error contributes at most
$N^{o(1)}K^{-1}\leq N^{o(1)}r$.

Use \eqref{loc:time-frostman} for that encounter interval and
\eqref{loc:conditional-cap} at each admissible time.  Their product
is at most $N^{o(1)}r^{s+a_{\mathrm{ang}}}$, as required.  This
argument is conditional probability followed by integration; it
does not assume that the encounter time and direction are independent.

\paragraph{Plate density and the filled support.}
The line-parameter measure is now $t$-Frostman.  We apply the planar
theorem to the neighbors' shadings, and use the resulting plate
density to fill the reference blocks.
Each neighbor index carries its original shading on the entire
current block.  It consists of
$KN^{-F(\kappa)+o(1)}$ cells, with the $s$-Frostman bound
\eqref{loc:time-frostman}.  Repetitions of a neighbor through different
sampled times merely add its line weight; they do not shorten this
shading.  The bounded projection preserves its cell count and
Frostman estimates.  Thus the weighted and repeated-line formulation
in Corollary~\ref{cor:weighted-furstenberg} applies to this probability
law and these shadings.  It gives at least
\begin{equation}
 N^{-F(\kappa)+o(1)}K^{1+\mu}
 \label{loc:planar-union-count}
\end{equation}
planar cells in their union.  More explicitly, one first fixes any
small theorem loss, includes all required cell and conditional mass
tests at that tolerance, and then lets the loss decrease.  The
strict inequality tested below is fixed before these choices.  This
is sufficient for the subpower notation in
\eqref{loc:planar-union-count}.

All the original neighboring segments lie in the plate from
Lemma~\ref{loc:short-block-plane}.  Its thickness is bounded, so a
projected cell receives only boundedly many relevant normal layers.
In particular \eqref{loc:planar-union-count} lower-bounds the number
of original union cells in this plate.  The plate has volume
comparable to $K^2$.  It therefore has original-union density at least
\begin{equation}
 \rho\geq N^{-F(\kappa)+o(1)}K^{\mu-1}.
 \label{loc:prefix-density}
\end{equation}
It contains the full reference segment, including times that were
not originally marked.  The same construction applies to every
retained reference index-block.

Fill all such reference blocks with every fine time bin, keeping
each original index and its inherited weight.  Write $E^+$ for the
new support.  Every point of $E^+$ lies in a plate with density
\eqref{loc:prefix-density} in a fixed enlargement of the original
support.  Lemma~\ref{lem:plate-nikodym} gives
\begin{equation}
 |E^+|\leq
 N^{2F(\kappa)+2\kappa(1-\mu)+o(1)}|E|.
 \label{loc:prefix-support-cost}
\end{equation}
The squared density loss in this formula is retained as a power of
$N$.

The number of incidences increases by
$N^{F(\kappa)+o(1)}$.  Indeed each retained occupied block had
$KN^{-F(\kappa)+o(1)}$ times and now has $K$.  The filled profile
is zero below $\kappa$ and equals $F(u)-F(\kappa)$ above it.
Its cost is $\Pi(1)-\Pi(\kappa)$, while its global raw clustering
parameter is at most $\Delta_{h_0}$.  The support estimate gives
\[
 m^+\geq
 mN^{-F(\kappa)-2\kappa(1-\mu)-o(1)}.
\]
The sharp hybrid upper bound gives
$m^+\leq mN^{-\Pi(\kappa)+o(1)}$.
These are incompatible with a strict reverse inequality to
\eqref{loc:prefix-necessary}.  This proves the lemma.
\end{proof}

\begin{lemma}[Positive flat delay]
\label{lem:flat-delay}
Every hybrid or exact equality with $H>0$ has a profile that is
identically zero on an interval $[0,\tau]$ with $\tau>0$.
\end{lemma}

\begin{proof}
Suppose that there is no such interval, and put
$e=\liminf_{u\downarrow0}F(u)/u$.  Then $F(u)>0$ for every $u>0$.
For any sufficiently small $\varepsilon>0$, choose a small $r_0>0$
so that $F(u)\geq(e-\varepsilon)u$ for $0\leq u\leq r_0$.
There are arbitrarily small $r\leq r_0$ where
$F(r)<(e+\varepsilon/2)r$.  The continuous function
$F(u)-(e+\varepsilon)u$ has a negative minimum on $[0,r]$.
Choose a minimizing point $\kappa>0$.  It is a record minimum
on $[0,\kappa]$, and hence
\begin{equation}
 F(\kappa)-F(u)\leq(e+\varepsilon)(\kappa-u),\qquad
 F(u)\geq\max(e-\varepsilon,0)u
 \quad(0\leq u\leq\kappa).
 \label{loc:record-prefix}
\end{equation}
We may take $\kappa$ as small as needed for
Lemma~\ref{lem:prefix-extension}.

At a typical fine incidence the neighboring mass is at least
$mN^{-o(1)}$ and lies within $O(K^{-1})$ of one plane.  Cover these
directions by $O(K)$ caps of radius $K^{-1}$ and use
\eqref{loc:cap-bound} together with \eqref{loc:record-prefix} and
\eqref{loc:cost-lower}.  This gives
$1\leq N^{o(1)}K^{1-H-p_*\max(e-\varepsilon,0)}$, hence
$H+p_*\max(e-\varepsilon,0)\leq1$, without first requiring a
positive time exponent.  Letting
$\varepsilon\downarrow0$ yields $H+p_*e\leq1$.
In particular $e<1$, and we may arrange
$s=1-e-\varepsilon>0$.  Use
\eqref{loc:record-prefix} in Lemma~\ref{lem:prefix-extension}, with
\[
 \gamma=\max(e-\varepsilon,0),\qquad
 t=1-e-\varepsilon+H+p_*\gamma.
\]

If $e=0$, choose $\varepsilon$ small compared with $H$.  Then
$s=1-\varepsilon$, $t=1+H-\varepsilon$, and
$\mu=1$.  Since $F(\kappa)>0$,
\[
 \Pi(\kappa)-F(\kappa)
       \geq(p_*-1)F(\kappa)>0,
\]
contradicting \eqref{loc:prefix-necessary}.  This argument does not
need a lower bound for $F(\kappa)/\kappa$.

Suppose $e>0$.  At $\varepsilon=0$ one has
\[
 s=1-e,\qquad t=1+H+(p_*-1)e>1,
\]
and therefore
\[
 2(1-\mu)=\max\{0,(2-p_*)e-H\}.
\]
The gap between $(p_*-1)e$ and the right-hand side is strictly
positive.  If the maximum is zero this is immediate; otherwise
the gap equals $(2p_*-3)e+H>0$.  By continuity, for sufficiently
small fixed $\varepsilon$,
\[
 (p_*-1)(e-\varepsilon)>2(1-\mu).
\]
But \eqref{loc:record-prefix} and
\eqref{loc:cost-lower} now give
\[
 \Pi(\kappa)-F(\kappa)
 \geq(p_*-1)F(\kappa)
 \geq(p_*-1)(e-\varepsilon)\kappa
 >2\kappa(1-\mu),
\]
again contradicting \eqref{loc:prefix-necessary}.  Both cases are
impossible, proving the flat delay.
\end{proof}

We have proved that an equality has an initial interval with zero
temporal deficit.  On any fixed scale in that interval, its shadings
have subpower density loss, so the weighted set estimate applies.
Together with the local bound, it forces the following plate shape.

\begin{lemma}[Saturated plates on a flat scale]
\label{loc:flat-saturation}
At every fixed positive scale exponent $\kappa$ with
$F|_{[0,\kappa]}=0$, an extensive regular restriction admits
disjoint indexed-block assignments to plates with widths
$1,K$, up to subpowers, where $K=N^\kappa$.  Each active plate
packet has raw index mass at least $mK$ in exponent.
\end{lemma}

\begin{proof}
Split into temporal $K$-blocks and bounded-overlap position charts
of diameter $CK$ fine cells containing complete block segments.
Each original index-block uses boundedly many charts.  The normalized
resolution is $K$, and the local shading density is $N^{-o(1)}$
because $F(\kappa)=0$.  Count charts separately.  Their total
incidence is extensive and their aggregate support is at most a
fixed multiple of the old support, so every extensive remainder
has aggregate raw multiplicity at least $m$ in exponent.

Fix a small $\varepsilon>0$.  Greedily remove and assign all
remaining local lines in a plank whenever its density
$n_W/(ab)$ is at least $mN^{-\varepsilon}$.  If an extensive
remainder survives, its clustering parameter is at most
$mN^{-\varepsilon}$ in every local chart.  The weighted set
estimate, Lemma~\ref{lem:weighted-set}, applies with local
resolution $K$ and arbitrarily small fixed theorem loss: choose
that loss below $\varepsilon/(2\kappa)$, then take $N$ large enough
for the subpower local shading density to meet its hypothesis.
It bounds the remainder's multiplicity by at most
$mN^{-\varepsilon/2}$ in exponent.  Uniformity over charts follows
by representative selection.  This contradicts its extensive
incidence and inherited support.  Thus the captured graph is
extensive, with disjoint index ownership on each block.

Combine its lower bound $mN^{-\varepsilon}ab$ with
\eqref{loc:mass-bound}, where $\Pi(\kappa)=0$.  It gives
\[
 N^{-\varepsilon}\leq
 N^{o(1)}K^{-H}a^{-q-H}b^H,
\]
or, in logarithmic coordinates,
\[
 (q+H)\log_N a+H\log_N(K/b)
       \leq\varepsilon+o(1).
\]
Both terms are nonnegative and $H>0$.  Let the capture tolerance
decrease to zero after each finite set of tests.  The captured
planks have $a=N^{o(1)}$, $b=KN^{o(1)}$, and their mass is at
least $mK N^{-o(1)}$.  Cleaning the assignment and chart partitions
preserves these lower bounds on occupied packets, as in
Proposition~\ref{prop:local-bound}.
\end{proof}


\subsection{Anisotropic normalization with inherited weights}

The next transformation uses alignment of the saturated plates.
It preserves their thin spatial coordinate while compressing time
and the long spatial coordinate by the angular scale $\theta$.
At fixed packet mass, the incidence count involves $K\theta$
time bins and the support bound involves $(K\theta)^2$ cells.
Their ratio gives a raw multiplicity gain of $\theta^{-1}$.  The lemma
also records the pulled-back plank bound and the remaining temporal profile.

\begin{lemma}[Anisotropic normalization]
\label{lem:anisotropic}
Suppose that on an extensive regular graph, over larger blocks of
$K_1=N^{\kappa_1}$ bins, there are disjoint indexed-block assignments
to saturated plates on a flat scale $K=N^\kappa\leq K_1$.  Their
widths are $1,K$ up to subpowers and their raw active index masses
are at least $mK$ in exponent.  Suppose the thin spatial axes of
the assigned plates along each original index and larger block
are within $\theta$ of one fixed axis for that index and block,
up to subpower factors, where
\[
 K^{-1}\leq\theta\leq1
 \quad\text{in exponent},\qquad
 \lambda=\log_N(1/\theta)\leq\kappa.
\]
There is an ensemble of bounded-chart configurations at resolution
\begin{equation}
 N'=K_1\theta
 \label{loc:anisotropic-resolution}
\end{equation}
whose aggregate raw multiplicity is at least $m/\theta$ in exponent.
Each derived index retains its original weight, and each original
index and larger block contributes to only subpower many charts.

Suppose also that all original active lines on a larger block obey
the local clustering bound $n_W\leexp C a^{x'}b^{y'}$ for
$1\leq a\leq b\leq K_1$, where $0\leq x'\leq y'$.  Then every
derived chart has full-time clustering parameter at most
\begin{equation}
 C\theta^{-y'}
 \label{loc:anisotropic-parameter}
\end{equation}
for those same shape exponents.
If $\kappa_1-\lambda>0$, its temporal profile, uniformly over the
ensemble, is
\begin{equation}
 G(v)=\frac{F(\lambda+(\kappa_1-\lambda)v)}
              {\kappa_1-\lambda},\qquad0\leq v\leq1.
 \label{loc:anisotropic-profile}
\end{equation}
Its cost in original $\log N$ units is $\Pi(\kappa_1)$; only the
initial flat interval $[0,\lambda]$ has been deleted.
\end{lemma}

\begin{proof}
\emph{The coordinate map and multiplicity gain.}
Replacing $\theta$ by $\min(\theta,1)$ changes it only by a
subpower factor, so we may take $0<\theta\le1$ in the geometric
construction.  First consider a chart centered on one affine line, with
normal--tangent coordinates $z_n,z_t$ and time $r\in[0,K_1]$, all
in old fine-cell units.  The chart bounds we will need are
$|z_n|\lesssim\theta K_1$, $|z_t|\lesssim K_1$, together with the
corresponding bounds on displacement over the larger block.
For such a chart use normalized coordinates
\begin{equation}
 s=\frac r{K_1},\qquad
 X=\frac{z_n}{\theta K_1},\qquad
 Y=\frac{z_t}{K_1},
 \label{loc:anisotropic-coordinates}
\end{equation}
at resolution $N'=K_1\theta$.  Equivalently, in the new fine-cell
units the coordinate map is
\begin{equation}
 (r,z_n,z_t)\longmapsto(\theta r,z_n,\theta z_t).
 \label{loc:anisotropic-fine-map}
\end{equation}
The new slopes and intercepts are bounded.  An old normal fine
cell remains a fine cell; old time and tangent cells are grouped
in intervals of length $1/\theta$.

Suppose a fine plate's axis differs from the chart normal by
$|\delta|\le\theta/10$, as the angular partition below will ensure.
Its thin condition is transformed from
$|\cos\delta\,z_n+\sin\delta\,z_t-c(r)|\lesssim N^{o(1)}$
to
\[
 |\cos\delta\,z_n'+(\sin\delta/\theta)z_t'-c'(r')|
       \lesssim N^{o(1)}.
\]
The coefficients have bounded or subpower size and the normal
coefficient is bounded below.  Hence it has subpower fine thickness
and tangential extent at most $N^{o(1)}K\theta$ in the new chart.
Fill its selected $K$-intervals.  Each now spans $K\theta$
new bins, up to rounding, and occupies at most
$N^{o(1)}(K\theta)^2$ cells.  If its retained index mass is at least
$mKN^{-o(1)}$, its incidence mass is at least
$mKN^{-o(1)}\cdot K\theta$.  Thus disjoint indexed-block
assignments with these retained masses
give aggregate raw multiplicity at least $m/\theta$ in exponent.
Packet support overlaps can only improve this lower bound.

\paragraph{Charts and preservation of the time tree.}
We now construct charts with the stated bounds, keeping each
original index and its selected time blocks together and retaining
the packet masses required by the cell count.
Use angular bins of width $\theta/10$ for the unoriented thin spatial
axis.  Since the axes used along any one original index and larger
block lie within subpower times $\theta$ of its fixed axis, that
index-block meets only subpower many bins.  In each angular bin
choose a normal--tangent frame.  Subdivide matrix space, expressed
over the larger interval, into parent cells of widths
$\theta K_1,K_1$ in normal and tangent position and also in the
normal and tangent displacement over that interval.  A fixed
matrix has only boundedly many adjacent grid parents in each bin.

An entire fine plate packet meets only subpower many such parents.
Indeed its thin axis belongs to the angular bin and differs from
the bin normal by at most $\theta/10$.  Its cross-sectional spread measured in that
normal is at most $N^{o(1)}(1+\theta K)$, and its tangential spread
is at most $N^{o(1)}K$.  Divide the endpoint differences by the
fine-block duration, then extrapolate over the larger interval.
The resulting spreads are at most
\[
 N^{o(1)}(K_1/K+\theta K_1)
       \leq N^{o(1)}\theta K_1
 \quad\text{in the normal coordinate},
\]
and $N^{o(1)}K_1$ in the tangent coordinate.  This uses
$\theta\geq K^{-1}$ in exponent.  The same bounds hold for
intercepts after subtracting a parent center.  Thus the asserted
parent routing follows from endpoint extrapolation, rather than
from the number of fine subblocks used by the index.

Within a chosen larger block, bin, and matrix parent, keep all
selected fine intervals of one original index together as a single
derived index.  Do not create a separate derived index for each
fine interval.  This preserves inherited weights and ensures that
a full-plank test in a chart counts a subset of original active
lines on the larger block.  Include packet subdivisions and parent
assignments in the finite cleaning.  Since every original
index-block has subpower routing, Corollary~\ref{cor:active-weight-cleaning}
shows that surviving packets retain active mass $mKN^{-o(1)}$.

The routing also preserves the time tree.  On every fixed logarithmic
subrange of the larger block, each derived time tree is contained in its original tree.  Its branching
exponents therefore cannot increase.  The chart ensemble retains
an extensive part of the weighted incidence, while its total
available original index-block mass has increased by only a
subpower.  After uniform regularization there can be no loss in
the total branching exponent.  The nonnegative losses on each
fixed subrange must consequently all be zero.  This is the
ensemble profile-preservation argument of
Lemma~\ref{lem:regularization}, with its routing hypothesis now
verified.  Thus the inherited prefix is $F|_{[0,\kappa_1]}$.
The charts and retained packets now meet the hypotheses of the first
calculation.  Applying \eqref{loc:anisotropic-fine-map} and filling
the selected blocks gives the asserted lower multiplicity $m/\theta$.

\paragraph{The plank parameter and temporal cost.}
For the clustering calculation, pull back a new full-time plank of
spatial widths $a,b$ by the inverse of
$(z_n,z_t)\mapsto(z_n,\theta z_t)$.  Its cross-section is a
parallelogram contained in a rotated rectangle of side lengths
$A\leq B$ with
\[
 AB\lesssim ab/\theta,\qquad B\lesssim b/\theta.
\]
For example, singular-value axes of the linear image of its
inscribed ellipse give this rectangle with fixed comparison
constants.  Its smallest width is bounded below by a constant
times $a$, and $B\lesssim b/\theta\leq K_1$; thus these are
permitted old tests after bounded coverings or clipping.  Since
$0\leq x'\leq y'$, the pulled-back mass is at most
\begin{align*}
 C A^{x'}B^{y'}
 &=C(AB)^{x'}B^{y'-x'}\\
 &\lesssim C(ab/\theta)^{x'}(b/\theta)^{y'-x'}
   =C\theta^{-y'}a^{x'}b^{y'}.
\end{align*}
Each chart index is an original index counted only once there,
so no extra weight factor enters this inequality.  This proves
\eqref{loc:anisotropic-parameter}.

Finally, the fine scale $1/\theta=N^\lambda$ lies in the flat
prefix, since $\lambda\leq\kappa$.  Filling the selected
$K$-blocks changes neither the limiting profile nor its cost.
Rebinning at $1/\theta$ deletes the initial flat prefix of length
$\lambda$ and translates the remaining logarithmic scale axis.
After dividing by the new total logarithmic length
$\kappa_1-\lambda$, this is precisely
\eqref{loc:anisotropic-profile}.  Its cost in original units is
$\Pi(\kappa_1)-\Pi(\lambda)=\Pi(\kappa_1)$.
When $\kappa_1-\lambda=0$ in the limit, the packet and clustering
bounds still have their asserted exponents; an upper estimate
can use total mass instead of a positive-resolution profile.
\end{proof}

\subsection{An equality with positive total deficit}

The two rescalings retain complementary parts of the time profile.
The isotropic cut keeps its branching above $K$ and fills the scales
below that block.  The anisotropic normalization keeps the prefix
through $K_1$ and removes only its initial flat range
$[0,\lambda]$.  We will use both on exact equalities in the next
section.  First we use the anisotropic operation in the hybrid case,
where the strict inequality $H>h_0$ gives a contradiction if the
plates can be aligned over the full time interval.

\begin{lemma}[Positive-deficit equality]
\label{lem:positive-deficit}
For $h_0\uparrow h>0$, the total deficits $\alpha=F(1)$ of hybrid
equalities are bounded below by a positive constant, uniformly
once $h_0$ is sufficiently close to $h$.  Consequently there is
an exact equality at $H=h_0=h$ whose limiting profile satisfies
$F(1)>0$.
\end{lemma}

\begin{proof}
We show that a sufficiently small fixed upper bound on $\alpha$
is incompatible with any such hybrid equality.  All fixed
tolerances in the next paragraphs will be chosen independently
of $H-h_0$.  Restrict to $h_0\geq h/2$, so that $H\geq h/2$.

We will capture an extensive incidence graph in nearly maximally elongated
global slabs, then align the saturated local plates with the fixed slab axis
assigned to each index. Anisotropic normalization will make the sharp hybrid
upper bound smaller than the retained raw multiplicity.

\paragraph{Fine-scale planarity and global slab capture.}
We first establish a fine-scale version of the multilinear
discard.  Choose a small constant $c_*>0$, depending only on the
positive lower bound $h/2$ and the multilinear input.  At a fine
cell use direction caps of radius $N^{-2c_*}$, declaring a cap
rich if its active mass is at least $mN^{-10c_*}$.  Mark a cell
bad if three rich caps have determinant at least $N^{-c_*}$.
The same multilinear count as in \eqref{loc:bad-balls}, now at
fine resolution, gives
\[
 \#\{\text{bad fine cells}\}
 \lesssim_\varepsilon
 N^{\varepsilon+C_3c_*}(n/m)^{3/2}
\]
for a fixed $C_3$.  Local bound \eqref{loc:mass-bound} at $K=1$
bounds active mass at a fine cell by $mN^{o(1)}$.  By
\eqref{loc:multilinear-slack}, the fraction of incidence in bad
cells is at most
\[
 N^{\varepsilon+C_3c_*+o(1)}
 \left(\frac{n}{k^2m}\right)^{1/2}
 \leq N^{\varepsilon+C_3c_*-h/4+o(1)}.
\]
Choose $c_*$ and then $\varepsilon$ so this is a negative fixed
power.  At every remaining cell all triples of directions in
rich caps have determinant $O(N^{-c_*})$.  The total mass in
light caps is at most $CmN^{-6c_*}$, since there are
$O(N^{4c_*})$ caps.  These choices are independent of any later
flat scale.

Choose a fixed $e_1>0$ sufficiently small compared with $c_*h$
and $h$, and require $\alpha$ to be smaller still, as well as
small enough for Lemma~\ref{lem:weighted-set} with target loss
$e_1/2$.  Greedily assign and remove all original full-time
indices in any global plank $S$ of widths $a,b$ with density
\begin{equation}
 n_S/(ab)\geq mN^{-e_1}.
 \label{loc:global-slab-capture}
\end{equation}
If an extensive remainder survives, its profile remains $F$,
so its shading density is $N^{-\alpha+o(1)}$.  All its plank
densities are at most $mN^{-e_1}$.  The weighted set estimate
then gives multiplicity at most $mN^{-e_1/2}$, contrary to
extensivity and the inherited support bound.  Thus the assigned
graph is extensive.  Each original index has exactly one global
assigned slab, and all of its retained times keep that label.

Compare \eqref{loc:global-slab-capture} with the original
$\Delta_{h_0}$ bound and the hybrid equality.  Put
$A=\log_N a$, $B=\log_N b$ just for this calculation.  Then
\begin{equation}
 (q+h_0)A+h_0(1-B)+(H-h_0)+\Pi(1)
       \leq e_1+o(1).
 \label{loc:global-slab-shape}
\end{equation}
In particular
\[
 a\leq N^{O(e_1/h_0)},\qquad
 b\geq N^{1-O(e_1/h_0)}.
\]
If \eqref{loc:global-slab-shape} has no feasible shape, the capture
already gives a contradiction; otherwise these width bounds hold
on every retained slab.

The assigned slab $S$ has at least $n_Sk$ weighted incidences and
support in $O(abN)$ cells.  Its average assigned multiplicity is
therefore at least $mN^{-e_1-\alpha-o(1)}$.  In each slab delete
cells with assigned multiplicity smaller by an inverse subpower
factor.  Their incidence is bounded by that factor times the
slab incidence, by the support count and
\eqref{loc:global-slab-capture}.  Summing over slabs and cleaning
the corresponding partitions gives the following lower bound at
every retained query incidence:
\begin{equation}
 \mu_S\geq mN^{-e_1-\alpha-o(1)}.
 \label{loc:slab-query-mass}
\end{equation}
This is a lower bound for the incidence from its own assigned
slab.  No independence from any later assignment is asserted.

\paragraph{Alignment with saturated local plates.}
The global slab assigned to an index has a fixed axis.  We next
compare it with the axes of the local saturated plates used by that
index.  By Lemma~\ref{lem:flat-delay}, this particular hybrid
equality has a positive flat delay.  Now choose a fixed
\[
 0<\xi\ll c_*
\]
smaller than that delay, and set $K=N^\xi$.
Lemma~\ref{loc:flat-saturation}, applied within the extensive
$S$-assigned graph, supplies saturated local plates $W$ of
widths $1,K$ up to subpowers and mass at least $mK$ in exponent.
Each such packet has $mK^2$ incidence in exponent and at most
$K^2N^{o(1)}$ support cells.  The analogous low-cell deletion
therefore gives, at its retained query incidences,
\begin{equation}
 \mu_W\geq mN^{-o(1)}.
 \label{loc:local-plate-query-mass}
\end{equation}
Retain \eqref{loc:slab-query-mass} as another partition test.
Both lower bounds thus hold on the same extensive central graph.
Fine multilinear bad cells can be discarded from it by their
power-small original incidence bound, followed by the same
subpower deficiency cleaning.

At such a query cell, remove directions in the global light caps
when choosing comparison directions.  Their mass
$CmN^{-6c_*}$ is negligible compared with both
\eqref{loc:slab-query-mass} and
\eqref{loc:local-plate-query-mass}, by the choices of $e_1$ and
$\alpha$.  Among the remaining directions in $S$ there are two
separated by at least $N^{-c_*/10}$.  Otherwise one cap of this
radius would contain all their mass, whereas
\eqref{loc:cap-bound} with $L=N^{c_*/10}$ bounds that cap mass
by $mN^{-Hc_*/10+o(1)}$.  Choose
$e_1+\alpha<Hc_*/20$ to contradict
\eqref{loc:slab-query-mass}.  Likewise there are two remaining
directions in $W$ separated by at least $K^{-1/2}$: the cap
bound with $L=K^{1/2}$ is $mK^{-H/2}N^{o(1)}$, which is smaller
than \eqref{loc:local-plate-query-mass} for this fixed $\xi>0$.

We justify how these two pairs compare the actual plate axes.
A plank with spatial thin axis $n_0$, thickness $a_0$, and
duration $K_0$ confines $n_0\cdot u$ to an interval of length
$O(a_0/K_0)$.  Its direction plane has normal proportional to
$(-v_0,n_0)$, where $v_0$ is the central normal slope.  Since
the original slopes are bounded, this normal has spatial
component bounded below after normalization.  Thus all directions
in $S$ are within
$N^{-1+O(e_1/h_0)}$ of its actual direction plane; all directions
in $W$ are within $K^{-1+o(1)}$ of its direction plane.

Let $v_1,v_2$ be the separated rich pair from $S$.
Every rich direction is within
\[
 O\!\left(\frac{N^{-c_*}}{|v_1\wedge v_2|}\right)=O(N^{-9c_*/10})
\]
of their span,
by the absence of a bad triple.  The span of $v_1,v_2$ is within
$N^{-1+O(e_1/h_0)+c_*/10}$ of the actual $S$ plane.
If $w_1,w_2$ is the separated rich pair from $W$, their span is
within $N^{-9c_*/10+\xi/2+o(1)}$ of the span of $v_1,v_2$,
and within $K^{-1/2+o(1)}$ of the actual $W$ plane.
Consequently the two actual direction planes differ by at most
\begin{equation}
 N^{-1+O(e_1/h_0)+c_*/10}
 +N^{-9c_*/10+\xi/2+o(1)}
 +N^{-\xi/2+o(1)}.
 \label{loc:axis-error}
\end{equation}
With the fixed choices above and $\xi\ll c_*$, this is smaller
than $\theta=N^{-\xi/4}$.  Taking the normalized spatial
components of the plane normals is a uniformly Lipschitz map
here.  Hence the thin spatial axes of $W$ and of the index's
assigned global slab $S$ differ, up to sign, by at most $C\theta$.

\paragraph{The hybrid contradiction.}
For each original index its $S$ axis is fixed over all time.
Thus the hypotheses of Lemma~\ref{lem:anisotropic} hold globally,
with $K_1=N$ and the finer flat scale $K=N^\xi$.
All previous discards were extensive or removed a power-small
incidence fraction.  The fixed losses $e_1$ and $\alpha$ were
used only to obtain the query lower bounds and identify the
axes; they are not losses in the retained incidence or packet
mass.  After subpower cleaning and routing, the lower raw
multiplicity supplied by that lemma is therefore $m/\theta$,
without an extra factor $N^{-e_1-\alpha}$.

We now use the strict gap between the hybrid exponents.  Take
$(x',y')=(d-1-h_0,1+h_0)$ in the anisotropic clustering bound; these
satisfy $0\leq x'\leq y'$.  The new resolution is
$N\theta$, and the temporal cost in original units is still
$\Pi(1)$.  The sharp hybrid upper estimate becomes
\begin{align}
 m_{\mathrm{new}}
 &\leexp\Delta_{h_0}\theta^{-(1+h_0)}
                  (N\theta)^H N^{\Pi(1)}\notag\\
 &\eqexp m\theta^{H-1-h_0}.
 \label{loc:hybrid-anisotropic-upper}
\end{align}
Its ratio to the lower bound $m/\theta$ is
$\theta^{H-h_0}$, a fixed negative-power saving in $N$ because
$H-h_0>0$ and $\xi>0$ are fixed for this hybrid equality.
This is a contradiction.  The possibly small value of
$\xi(H-h_0)$ causes no problem: that hybrid and its flat scale
are fixed before its realizing resolution tends to infinity.

\paragraph{Passage to an exact equality.}
The choices of $c_*,e_1$ and the upper bound on $\alpha$ depended
only on $h,p,q$ and the external theorem tolerances, not on
$h-h_0$.  This proves a uniform positive lower bound for
$\alpha$ as $h_0\uparrow h$.
Finally, for every configuration,
\[
 N^{-(h-h_0)}\Delta_{h_0}\leq\Delta_h\leq\Delta_{h_0},
\]
because $1\leq b/a\leq N$ in each plank test.
Thus the hybrid equality exponents tend to the exact critical
one.  Apply the profile compactness and equality diagonal from
Lemma~\ref{lem:hybrid-equality}.  A downward jump of the convex
cost would improve the critical score and contradict the sharp
upper bound, so it cannot spoil equality.  Uniform convergence
of the profiles preserves their positive endpoint lower bound.
The resulting exact equality has $F(1)>0$, as claimed.
\end{proof}
